\chapter{Forthcoming Course Modules}
\label{ch:forthcoming}

\begin{plainlanguage}
Edition 0.1 ends after introductory SQL and the cumulative lab. The modules in
this chapter are a transparent roadmap, not claims of finished instruction.
Each entry names the intended learning work and the project evidence that a
later edition should use.
\end{plainlanguage}

\ForthcomingChapter
  {SQL Joins and Analytical Grain}
  {Join the OULAD relations without changing the intended unit of analysis;
   reason about one-to-one, one-to-many, and many-to-many paths; detect row
   multiplication; and validate pre-aggregation.}
  {Construct an attempt-level assessment summary and an attempt-level activity
   summary, then join them only after key tests pass.}
  {\projectfile{sql/02_staging.sql}, \projectfile{sql/03_core.sql}, and
   \projectfile{sql/gallery/}.}

\ForthcomingChapter
  {Intermediate and Advanced SQL}
  {Use common table expressions, window functions, conditional aggregation,
   date-relative frames, query plans, views, and materialized views while
   preserving analytical meaning.}
  {Build presentation-level rankings and cumulative engagement summaries, then
   compare equivalent query plans.}
  {\projectfile{sql/gallery/}, \projectfile{sql/06_features.sql}, and existing
   reporting queries.}

\ForthcomingChapter
  {Python and pandas}
  {Load database extracts safely, inspect dtypes and keys, write vectorized
   transformations, test reusable functions, and compare pandas operations with
   their SQL equivalents.}
  {Reproduce one validated SQL summary in pandas and reconcile every count.}
  {\projectfile{src/learning_analytics/config.py},
   \projectfile{src/learning_analytics/analysis.py}, and the test suite.}

\ForthcomingChapter
  {Exploratory Data Analysis}
  {Separate data checks from exploration; examine distributions, missingness,
   time, and subgroup composition; and maintain a question-and-evidence log.}
  {Create an attempt-level EDA notebook or script whose reported totals reconcile
   to the quality outputs.}
  {\begin{itemize}
     \item \projectfile{reports/tables/source_file_profile.csv}
     \item \projectfile{reports/tables/sql_quality_results.csv}
     \item Existing analysis code
   \end{itemize}}

\ForthcomingChapter
  {Statistics}
  {Use sampling language, uncertainty intervals, effect sizes, and carefully
   bounded comparisons; distinguish statistical association from practical or
   causal importance.}
  {Estimate and interpret uncertainty for selected descriptive proportions
   without presenting the source as a random sample of all learners.}
  {Outcome summaries, the OULAD documentation, and later analysis outputs.}

\ForthcomingChapter
  {Visualization}
  {Choose visual encodings that match data type and question, expose
   denominators, handle missing categories, and test charts for accessibility
   and interpretive honesty.}
  {Rebuild the outcome distribution and one time-based view with a written
   chart-validation checklist.}
  {\projectfile{reports/figures/} and
   \projectfile{src/learning_analytics/analysis.py}.}

\ForthcomingChapter
  {Engagement Feature Engineering}
  {Define activity windows, aggregate clicks and active days, distinguish site
   breadth from volume, and test feature grain.}
  {Build attempt-level engagement features for explicitly bounded observation
   windows and reconcile them to source activity.}
  {\projectfile{sql/06_features.sql} and
   \projectfile{src/learning_analytics/features.py}.}

\ForthcomingChapter
  {Assessment Feature Engineering}
  {Construct submission, score, deadline, and completion features while
   preserving missingness and avoiding duplicated assessments.}
  {Produce an attempt-level assessment feature table with key, range, and
   denominator tests.}
  {\projectfile{sql/06_features.sql}, assessment source tables, and feature
   code.}

\ForthcomingChapter
  {Time-aware Analytical Design}
  {Define an observation date, prediction horizon, eligibility rule, and
   outcome window; align every feature to information available at that time.}
  {Draw and implement a timeline for one early-support analysis.}
  {Relative-day fields, registration timing, VLE activity, and assessment due
   dates.}

\ForthcomingChapter
  {Leakage Prevention}
  {Recognize target, temporal, aggregation, and split leakage; audit feature
   provenance; and write automated leakage checks.}
  {Classify candidate features by availability and reject those containing
   post-cutoff or outcome-derived information.}
  {\projectfile{docs/methodology.md}, feature SQL, and modeling code.}

\ForthcomingChapter
  {Predictive Modeling}
  {Establish a baseline, define a target and eligible cohort, fit interpretable
   models, and keep preprocessing inside the training workflow.}
  {Compare a naive baseline with a regularized classifier using a time-aware
   split.}
  {\projectfile{src/learning_analytics/modeling.py} and model configuration.}

\ForthcomingChapter
  {Model Evaluation}
  {Choose metrics for class imbalance and decision context; evaluate
   discrimination, calibration, thresholds, and subgroup behavior.}
  {Create a model-comparison table whose metrics are tied to a stated use case
   rather than a single winning score.}
  {Saved model metrics, prediction outputs, and
   \projectfile{docs/forecast-evaluation.md}.}

\ForthcomingChapter
  {Error Analysis and Interpretation}
  {Inspect false positives and false negatives, compare errors across relevant
   groups and presentations, and separate model behavior from causal stories.}
  {Build a reproducible error slice and write an interpretation with competing
   explanations.}
  {Prediction outputs, attempt attributes, and model documentation.}

\ForthcomingChapter
  {Responsible Educational Analytics}
  {Analyze privacy, consent, fairness, intervention burden, human oversight,
   and the consequences of acting on uncertain predictions.}
  {Write a pre-deployment review that includes a non-deployment option and an
   appeals path.}
  {\projectfile{docs/forecast-evaluation.md}, dataset documentation, and published OULAD
   context.}

\ForthcomingChapter
  {Reproducibility and Software Quality}
  {Package configuration, tests, data contracts, deterministic builds, logs,
   and documentation into an auditable analytical product.}
  {Rebuild the project from a clean environment and produce a provenance
   manifest for every reported artifact.}
  {\projectfile{Makefile}, \projectfile{pyproject.toml}, scripts, SQL entry
   points, and tests.}

\ForthcomingChapter
  {Final Capstone}
  {Integrate question design, relational data work, SQL, Python, statistics,
   visualization, time-aware modeling, careful interpretation, and
   reproducibility.}
  {Deliver an evidence-backed analytical report, reproducible code, quality
   record, model documentation where applicable, and a short oral defense.}
  {The complete learning-analytics repository and all completed course modules.}

\section{What the roadmap does not imply}

The presence of project scripts or saved artifacts does not mean the teaching
modules above are complete. A future edition must still explain each idea from
first principles, connect it to executable work, provide exercises and guided
solutions, and verify every reported result. Until then, these entries remain
explicitly forthcoming.
